\documentclass{report}

\usepackage{amsmath,amssymb,amsthm}
\usepackage{hyperref}
\usepackage[capitalize]{cleveref}

% Theorem environments
\newtheorem{theorem}{Theorem}
\newtheorem{definition}{Definition}
\newtheorem{axiom}{Axiom}
\newtheorem{corollary}{Corollary}
\newtheorem{remark}{Remark}

% Custom macros for leanblueprint
\providecommand{\home}[1]{}
\providecommand{\github}[1]{}
\providecommand{\dochome}[1]{}
\providecommand{\lean}[1]{}
\providecommand{\leanok}{}
\providecommand{\uses}[1]{}

\title{ISAR: A Quotient-Mediated Semantic Kernel for Closed Dialects}
\author{Fabian Schindler}
\date{\today}

\begin{document}
\maketitle

\begin{abstract}
This paper presents the ISAR framework, a closed combinatory substrate designed to unify disjoint computational models under a single algebraic setting. We formally prove that the quotient \texttt{InvariantLayer} is the view-independent invariant layer of the substrate. Furthermore, we establish that the representation process for any admissible closed dialect yields an $\text{Encoder} \to \text{Kernel} \to \text{Quotient} \to \text{Decoder}$ factorization, showing that different formalisms and programming languages function as decoders over the same invariant structure.
\end{abstract}

\input{content_paper_a.tex}

\end{document}
